% ============================================================
% CLASS CARD POOLS — ARCHETYPES (first draft)
% ------------------------------------------------------------
% AUTHOR SPEC: each class offers multiple POOLS, each pool an
% archetype characters of that class can fall into.
%
% STRUCTURE: 10 classes, 32 archetype pools. Every archetype
% gets: concept, playstyle, and 3–4 SAMPLE cards spanning early
% levels plus one signature — samples establish each pool's
% grammar; full pools (15–20 cards each) are follow-up work.
%
% [PROPOSALS FLAGGED FOR AUTHOR SIGN-OFF]
%   C1. Pool rule: at creation, choose ONE archetype; your 10
%       class cards come from its pool. From level 2 on, you may
%       buy from sibling archetypes of your class at +1 card
%       level (cross-training). Species/faction pools unaffected.
%   C2. Card line format: Name — Cost (Momentum) — Level — tags
%       — effect. "Usually 1" cost holds; big effects cost 2 or
%       carry One-Shot.
%   C3. New tags introduced: Reaction (playable during the enemy
%       side's activation at listed cost, no Momentum needed),
%       Troop (puts a unit in play), Stance (persists until you
%       play another Stance), Consumable (Sawbones/Tinkerer
%       items — may be handed to allies before a scene).
%   C4. Troop template (Captain): a Troop is a side member with
%       the statline printed on its card; its actions cost side
%       Momentum (1) like anyone's; it drops when its HP does;
%       troops re-muster between encounters unless the card says
%       One-Shot.
%   C5. Class name placeholder: the Monastic Orders class is
%       written here as "the Monk" pending the author's name.
%   C6. Shapeshifter FORMS come from the class's level-budget
%       feature (per author spec), not from cards; the pool cards
%       below are techniques laid over the forms.
%   C7. Numbers (damage, HP, dice) are tuning placeholders.
% ============================================================

\chapter{Class Card Pools}

\section{How Pools Work}

Every class teaches more than any one member learns. A class's full repertoire is divided into \textbf{archetype pools} --- distinct traditions, trainings, or expressions of the class --- and the archetype you choose at creation is the second-most defining choice you make after the class itself. [PROPOSAL C1] Your 10 starting class cards come from your archetype's pool; from level 2 onward you may cross-train, buying sibling-archetype cards at one card level higher than printed --- the techniques travel, but they were not built for your hands.

An archetype is not a subclass wall; it is a center of gravity. Two Zealots of different pools are recognizably the same calling and unmistakably different people, and the frontier knows the difference on sight --- which is the point. The archetypes below are the ones the setting's institutions actually produce, and each pool's flavor is that institution's fingerprint.

\medskip\hrule\medskip

\section{Zealot \textnormal{(Witnesses of the Lion)}}

The Lions' martial calling: close the distance, hold the line, and do not stop seeing.

\subsection{The Line-Holder}
\textit{The formation's anchor --- the one who is not moving, so that others don't have to be brave.} Playstyle: defense, protection, punishing enemies who ignore you.
\begin{itemize}
    \item \textbf{Hold Here} --- 1 --- L1 --- Stance --- While you don't move, allies within 2 tiles gain +1 Defend and cannot be forced to move.
    \item \textbf{You Face Me} --- 1 --- L1 --- One enemy within 3 tiles must include you in its next attack or lose 1 die from it.
    \item \textbf{The Door Does Not Open} --- 1 --- L3 --- Reaction --- When an enemy tries to pass within your reach, stop its movement and Strike it.
    \item \textbf{Redoubt} --- 2 --- L5 --- One-Shot --- Until your side's next surge, attacks against allies within 2 tiles of you target you instead, and you gain +2 Armor.
\end{itemize}

\subsection{The Intercessor}
\textit{The formation's honest voice --- prayer as logistics, aid arriving because it was asked for truthfully.} Playstyle: support, healing, buying allies' successes.
\begin{itemize}
    \item \textbf{Spoken Plainly} --- 1 --- L1 --- An ally within 5 tiles heals 2 HP or clears one fear or despair effect; you must say aloud, in character, why they will hold.
    \item \textbf{Borrowed Light} --- 1 --- L2 --- An ally's next roll this round gains +2 dice; if they tell the table what they're afraid of, +3.
    \item \textbf{The Vapor Carries It} --- 1 --- L4 --- One-Shot --- Ask the GM one true question about the present danger; the answer arrives as vision.
    \item \textbf{Answered} --- 2 --- L6 --- One-Shot --- Celestial aid, briefly: choose one --- all allies heal 3 HP; or one enemy of the Shadow takes hits equal to your Religion.
\end{itemize}

\subsection{The Lamplighter}
\textit{The burning end of witness --- the one sent where the dark is thickest, to make it regret the attention.} Playstyle: striker against the Shadow's works; light as a weapon.
\begin{itemize}
    \item \textbf{Lamp Oil and Iron} --- 1 --- L1 --- Your Strikes and Shoots this round ignore the Armor of undead and Shadow-aligned creatures.
    \item \textbf{Cast Light} --- 1 --- L1 --- Illuminate 5 tiles: hidden things are revealed; things that hate light lose 1 Defend there.
    \item \textbf{I See You} --- 1 --- L3 --- Name one enemy you can see: your side's attacks against it gain +1 die until it drops. It knows.
    \item \textbf{Dawn, Early} --- 2 --- L6 --- One-Shot --- A burst of revealing light: every enemy within 3 tiles takes 2 hits (Shadow-aligned: 4) and everything concealed is not.
\end{itemize}

\medskip\hrule\medskip

\section{Circuit Rider \textnormal{(Witnesses of the Lion)}}

The Lions' traveling ministry: the one who comes back, remembers, and carries.

\subsection{The Shepherd}
\textit{The circuit's healer --- of bodies first, and of the things bodies carry.} Playstyle: sustained healing, condition-clearing, keeping a side on its feet.
\begin{itemize}
    \item \textbf{Seen and Tended} --- 1 --- L1 --- Heal an adjacent ally 3 HP; if you've named a true thing about them this scene, 4.
    \item \textbf{Get Up. You Held.} --- 1 --- L2 --- Reaction --- An ally dropping to 0 HP drops to 1 instead.
    \item \textbf{Circuit Kit} --- 1 --- L3 --- Consumable ×3 --- Hand out before scenes: the bearer may spend it to heal 2 HP or reroll a save.
    \item \textbf{No One Dies Tonight} --- 2 --- L6 --- One-Shot --- Until end of scene, allies within 5 tiles cannot be reduced below 1 HP by the first source that would do it each round.
\end{itemize}

\subsection{The Witness}
\textit{The circuit's eyes --- the Rider who notices what changed, and says so.} Playstyle: information, exposure, denying enemies their deceptions.
\begin{itemize}
    \item \textbf{What Changed Here} --- 1 --- L1 --- Ask the GM: what in this scene is different from how it should be? True answer.
    \item \textbf{Say It Plain} --- 1 --- L2 --- One creature must answer your next question truthfully or visibly, unmistakably refuse.
    \item \textbf{The Account Is Kept} --- 1 --- L3 --- Stance --- Enemies within 5 tiles cannot gain benefits from concealment, disguise, or lies while you watch.
    \item \textbf{Enter It in the Book} --- 2 --- L5 --- One-Shot --- Name an enemy: for the rest of the scene, your whole side knows its position, and it cannot surprise anyone.
\end{itemize}

\subsection{The Long Rider}
\textit{The circuit itself, made flesh --- distance means less to them than it should.} Playstyle: mobility, positioning, logistics; the side arrives where it's needed.
\begin{itemize}
    \item \textbf{Road Legs} --- 1 --- L1 --- Stance --- You and adjacent allies gain +1 Move; difficult ground doesn't slow you.
    \item \textbf{I Know a Way} --- 1 --- L2 --- Move an ally within 5 tiles up to their Move (a shortcut, a boost, a shouted route).
    \item \textbf{The Lantern at Dusk} --- 1 --- L4 --- Until end of scene, allies who end their move adjacent to you heal 1 HP and clear fatigue.
    \item \textbf{Rider on the Ridge} --- 2 --- L6 --- One-Shot --- Help arrives: an NPC ally you plausibly know (GM's choice from your circuit) enters the scene's edge, favorably disposed.
\end{itemize}

\medskip\hrule\medskip

\section{Arcanist \textnormal{(Academic Enclaves)}}

Erupted power, stewarded: the mass, the book, the protocols.

\subsection{The Ward-Wright}
\textit{The useful Enclave, in one person --- containment as a combat art.} Playstyle: barriers, zones, control; the battlefield does what you chalk.
\begin{itemize}
    \item \textbf{Chalk Line} --- 1 --- L1 --- Draw a 3-tile line: enemies crossing it lose 1 die from their next action.
    \item \textbf{Warding Circle} --- 1 --- L2 --- A 1-tile circle: allies inside gain +1 Defend and +1 die on saves; lasts until you move.
    \item \textbf{Containment} --- 2 --- L4 --- Box one enemy in force (Defend 2, HP 4 walls); it batters or bypasses, but not this round.
    \item \textbf{The Vellhame Protocol} --- 2 --- L6 --- One-Shot --- Suppress all supernatural effects in a 3-tile radius for one round --- including your side's. It exists for a reason.
\end{itemize}

\subsection{The Open Channel}
\textit{The eruption, ridden --- output the protocols call inadvisable and the moment calls necessary.} Playstyle: blaster; push-your-luck damage that spends your own body.
\begin{itemize}
    \item \textbf{Arc} --- 1 --- L1 --- 3 hits to one target within 5 tiles; you may take 1 damage to make it 4.
    \item \textbf{Discharge} --- 1 --- L2 --- 2 hits to every enemy in a 2-tile cone; take 1 damage per extra tile of reach you add.
    \item \textbf{Overdraw} --- 1 --- L4 --- Stance --- Your Channel cards deal +1 hit; you take 1 damage at the start of each of your side's surges while this holds.
    \item \textbf{The Cascade, Glimpsed} --- 2 --- L7 --- One-Shot --- 6 hits, 3-tile radius, centered on you; allies included unless warded; the GM notes, quietly, that it felt easy.
\end{itemize}

\subsection{The Book-Bound}
\textit{The inheritance, leaned on --- forty dead masters between two chained covers, consulted mid-fight.} Playstyle: versatility and deck manipulation; the answer is in the book.
\begin{itemize}
    \item \textbf{Consult the Names} --- 1 --- L1 --- Look at the top 3 cards of your holstered pile; load one into your hand (discard one from hand if over 5).
    \item \textbf{A Dead Master's Trick} --- 1 --- L2 --- Copy the printed effect of any card in your expended pile, at +1 cost.
    \item \textbf{The Book Remembers} --- 1 --- L3 --- Reaction --- When an enemy uses a supernatural ability, ask the GM what it is and one weakness; the book has seen it before.
    \item \textbf{Forty Hands on the Pen} --- 2 --- L6 --- One-Shot --- Play up to 3 cards from your expended pile this surge, paying their costs; the book is briefly, unmistakably awake.
\end{itemize}

\medskip\hrule\medskip

\section{Dynast \textnormal{(Noble Houses)}}

The blood, expressed: every archetype is a lineage-pattern, and the great Houses are only the famous examples of each.

\subsection{The Thirsting Line \textnormal{(Corvara-pattern)}}
\textit{Speed past watching comfortably, and a debt the body collects.} Playstyle: skirmisher; hit-and-gone; blood as resource.
\begin{itemize}
    \item \textbf{Arrival Without Transit} --- 1 --- L1 --- Move your full Move and Strike; if the target hadn't seen you this round, +2 dice.
    \item \textbf{The Taste} --- 1 --- L2 --- After damaging a living creature in melee: heal 2 HP and learn one true physical fact about it (health, poison, lineage, expression).
    \item \textbf{Night's Own} --- 1 --- L3 --- Stance --- In dim light or darkness: +1 Move, +1 Defend, and you see perfectly.
    \item \textbf{The Red Ledger} --- 2 --- L6 --- One-Shot --- Against a creature you have tasted: this Strike ignores Defend, and you know where it is for the rest of the session.
\end{itemize}

\subsection{The Grown Line \textnormal{(Ostrogar-pattern)}}
\textit{Living armament --- plate that thickens where you were struck.} Playstyle: juggernaut; grows harder as the fight goes on.
\begin{itemize}
    \item \textbf{Blade-Ridge} --- 1 --- L1 --- Stance --- Your unarmed Strikes use grown edges: +1 die and lethal.
    \item \textbf{Answer in Bone} --- 1 --- L2 --- Reaction --- When an attack damages you, gain +1 Armor until end of scene (max +3); the plate remembers.
    \item \textbf{Horn and Weight} --- 1 --- L3 --- Move up to your Move in a line and Strike; on a hit, push the target 2 tiles and knock it down.
    \item \textbf{The Seat, Refused} --- 2 --- L7 --- One-Shot --- Until end of scene you cannot be moved, knocked down, or reduced below 1 HP by attacks; you also cannot move. A wall that remembers why it stands.
\end{itemize}

\subsection{The Living Iron \textnormal{(Ferrant-pattern)}}
\textit{Flesh of frame-metal --- a body that banks the charge.} Playstyle: durable battery; stored vim spent as violence or utility.
\begin{itemize}
    \item \textbf{Take the Charge} --- 1 --- L1 --- Absorb a vim source (a Heart, a jolt-line, an ally's cell): bank 3 \textit{ebbs} on this card's line (max 6).
    \item \textbf{Discharge Through the Hands} --- 1 --- L1 --- Spend banked ebbs: 1 hit per ebb to a creature you're touching, or power any device for a scene.
    \item \textbf{Ringing Bones} --- 1 --- L3 --- Stance --- +1 Armor; you are immune to vim/electrical harm and bank 1 ebb whenever it's tried.
    \item \textbf{Kin to the Herd} --- 2 --- L5 --- One-Shot --- A spagyric frame or chassis within sight treats you as its handler for the scene; it gentles, obeys, or carries.
\end{itemize}

\subsection{The Drowned Line \textnormal{(Ophara-pattern)}}
\textit{Scale, claw, and blood that fights back.} Playstyle: controller-bruiser; punishes attackers; owns the water.
\begin{itemize}
    \item \textbf{Acid in the Veins} --- 1 --- L1 --- Stance --- Creatures that hit you in melee take 1 hit; your dead body is a House matter, and everyone can tell.
    \item \textbf{The Drowned Court's Grip} --- 1 --- L2 --- Grapple with claws: +2 dice, and a grappled foe takes 1 acid hit per round held.
    \item \textbf{Regrow} --- 1 --- L3 --- Heal 3 HP; clear one lingering injury; over downtime, regrow what was lost.
    \item \textbf{Take It Under} --- 2 --- L6 --- One-Shot --- In or beside water: drag one enemy your size or smaller beneath; it spends its side's next 2 Momentum surfacing or begins to drown.
\end{itemize}

\medskip\hrule\medskip

\section{The Monk \textnormal{(Monastic Orders)} \textnormal{[name pending --- PROPOSAL C5]}}

The cultivated will: the weapon that cannot lie, and the two limbs of one way.

\subsection{The Edge}
\textit{The outward way --- what the blade can cut, climbed rung by rung.} Playstyle: striker whose weapon ignores what armor thinks matters.
\begin{itemize}
    \item \textbf{First Word} --- 1 --- L1 --- Stance --- Manifest your blade: Strikes with it use Nimbleness or Athletics dice (your choice) and are lethal.
    \item \textbf{Cut the Material} --- 1 --- L2 --- Your blade-Strikes this round ignore 2 Armor and can sever objects (chains, bars, timber) outright.
    \item \textbf{Cut the Working} --- 1 --- L4 --- Strike a spell, ward, or supernatural effect within reach: on 2 hits, it ends.
    \item \textbf{The Opened Way} --- 2 --- L6 --- One-Shot --- Cut a door in distance: you and adjacent allies step up to 10 tiles, through nothing, arriving instantly. The cut declines to go through anywhere.
\end{itemize}

\subsection{The Eye}
\textit{The inward way --- the still mind, and other minds, approached lawfully.} Playstyle: control and insight; the fight ends because someone understood it.
\begin{itemize}
    \item \textbf{Still Water} --- 1 --- L1 --- Stance --- +2 dice on saves against fear, deception, and mind-effects; you cannot be surprised.
    \item \textbf{Read the Stance} --- 1 --- L1 --- One creature you observe: learn its intent, its fear, or its next action (GM's honest answer).
    \item \textbf{Speech Without Voice} --- 1 --- L2 --- Until end of scene, your side converses silently at any distance in sight; coordination cards and aid cost allies 1 less die of penalty from range or noise.
    \item \textbf{The Closed Door, Held Open} --- 2 --- L7 --- One-Shot --- With consent, enter a willing mind: heal one mental affliction, share one memory truly, or stand with them against something that got in. Never without consent. The examiners check.
\end{itemize}

\subsection{The Walker}
\textit{The road as cell --- the wandering master, half festival, half inspection.} Playstyle: mobile support; teaches mid-fight; strong anywhere, supreme nowhere.
\begin{itemize}
    \item \textbf{Traveling Form} --- 1 --- L1 --- Stance --- +1 Move; opportunity Strikes against you lose 1 die; your blade manifests and dismisses freely.
    \item \textbf{A Correction, Offered} --- 1 --- L2 --- An ally within 3 tiles rerolls up to 2 dice of a roll they can see you watching.
    \item \textbf{The Gate Custom} --- 1 --- L3 --- Once per scene per creature: an enemy that accepts your offered pause stops fighting you until harmed; you both know it's real.
    \item \textbf{Learn What We Don't Know} --- 2 --- L5 --- One-Shot --- Observe any technique used this scene (card, deed, or ability): perform it once, now, at its printed cost.
\end{itemize}

\medskip\hrule\medskip

\section{Tinkerer \textnormal{(Landed Guilds)}}

The Procedures, carried in a satchel: your archetype is your hall of training.

\subsection{The Coilwright}
\textit{Energetics --- vim charged, kept, and delivered on schedule.} Playstyle: zone damage, batteries, powering the party's world.
\begin{itemize}
    \item \textbf{Jolt} --- 1 --- L1 --- 2 hits to a target within 3 tiles; constructs, frames, and Living Iron take none and bank it instead.
    \item \textbf{Storage Cell} --- 1 --- L1 --- Consumable ×2 --- A charged cell: the bearer may spend it for +1 Momentum for the side (their surge only) or to power a device for a scene.
    \item \textbf{Ground the Line} --- 1 --- L3 --- Reaction --- Redirect one vim, lightning, or Heart-driven effect targeting anything within 3 tiles into the earth. Bank 1 ebb; bill later.
    \item \textbf{The Dawn Clause} --- 2 --- L6 --- One-Shot --- At the procedurally correct moment (GM agrees the sky is right): your next Coilwright card triples its numbers. Nobody knows why. It replicates.
\end{itemize}

\subsection{The Crucible-Trained}
\textit{Alchemy --- matter argued into new opinions.} Playstyle: grenades, acids, transmutation; the environment is ammunition.
\begin{itemize}
    \item \textbf{Flask of Opinion} --- 1 --- L1 --- Thrown, 5 tiles: 2 hits in a 1-tile burst, and the ground is briefly (choose) burning, slick, or fogged.
    \item \textbf{Solvent} --- 1 --- L2 --- Destroy up to a door's worth of material, or strip 2 Armor from a target until repaired.
    \item \textbf{Alloy on Demand} --- 1 --- L3 --- Reforge or harden: an ally's weapon gains +1 die, or their Armor +1, until end of scene.
    \item \textbf{Orichalcum Sliver} --- 2 --- L6 --- One-Shot --- Overcharge a weapon with the crown metal: its next attack adds hits equal to your Study. The sliver is spent, and the Crucible will ask where you got it.
\end{itemize}

\subsection{The Frame-Wright}
\textit{Spagyrics --- the living machine, grown, gentled, and yours.} Playstyle: companion archetype; a small chassis fights beside you.
\begin{itemize}
    \item \textbf{The Frame} --- 1 --- L1 --- Troop --- Deploy your chassis: HP 5, Strike 3, Defend 1, Move 2, carries loads. It re-musters between encounters; it mourns if you don't.
    \item \textbf{Grown for This} --- 1 --- L2 --- Stance --- Choose the frame's graft for the scene: +1 Strike (goring horn), +1 Defend (plate), or +1 Move (long limbs).
    \item \textbf{Good Beast} --- 1 --- L2 --- Reaction --- The frame interposes: an attack on you or an adjacent ally hits the frame instead.
    \item \textbf{The Yard's Best Line} --- 2 --- L6 --- One-Shot --- The frame exceeds procedure: this scene it acts on 0-Momentum once per surge, and something in it is clearly, briefly, \textit{choosing} to. File a report. Or don't.
\end{itemize}

\medskip\hrule\medskip

\section{Sawbones \textnormal{(Landed Guilds)}}

The body, documented: healing with a case number.

\subsection{The Red Thread}
\textit{Chirurgeon-trained --- surgery under fire, anatomy as a map.} Playstyle: the party's true healer; precision offense in a pinch.
\begin{itemize}
    \item \textbf{Field Surgery} --- 1 --- L1 --- Heal an adjacent ally 3 HP and clear one bleed, poison, or injury; +1 HP if you narrate the procedure.
    \item \textbf{I Know Where It Hurts} --- 1 --- L2 --- Your Strike targets anatomy: on any hit, the target also loses 1 die from its next action.
    \item \textbf{Triage} --- 1 --- L3 --- Heal 2 HP split among up to 3 allies within 3 tiles; the case library note is free.
    \item \textbf{Not On My Table} --- 2 --- L6 --- One-Shot --- An ally at 0 HP or dead this scene is stabilized at 1 HP; the injury annex gains an entry with your name on it.
\end{itemize}

\subsection{The Preparatory}
\textit{Green Chain--trained --- tinctures, stims, and poultices, brewed ahead of need.} Playstyle: pre-battle buffer; consumables as party economy.
\begin{itemize}
    \item \textbf{Tincture Kit} --- 1 --- L1 --- Consumable ×3 --- Distribute before scenes: spend to heal 2 HP or clear fatigue.
    \item \textbf{Vigor Draught} --- 1 --- L2 --- Consumable ×2 --- Spend for +2 dice on one physical roll; the crash afterward is a role-playing opportunity.
    \item \textbf{Smelling Salts \& Worse} --- 1 --- L3 --- Reaction --- An ally dropping to 0 HP stays at 1 and acts normally this round; they owe you a conversation.
    \item \textbf{The Season's Brew} --- 2 --- L5 --- One-Shot --- Before a scene begins, every ally drinks: +1 die on all saves and +2 max HP for the scene.
\end{itemize}

\subsection{The Toxicologist}
\textit{Crucible-trained --- the dose makes the poison, and you did the math.} Playstyle: debuffer; the enemy's numbers quietly stop working.
\begin{itemize}
    \item \textbf{Coated Edge} --- 1 --- L1 --- An ally's weapon (or yours) is dosed: its next 2 hits also cost the target 1 die from its next action.
    \item \textbf{Know the Antidote} --- 1 --- L2 --- Reaction --- Cancel one poison, venom, or toxin effect on anyone within reach; identify its source.
    \item \textbf{The Slow Argument} --- 1 --- L4 --- Touch or hit: the target's dice pools drop by 1 at the start of each of its surges (max --3) until treated.
    \item \textbf{The Numbers Person} --- 2 --- L6 --- One-Shot --- Declare a creature type you've studied this session: your side's attacks against it gain +1 die and ignore 1 Armor for the scene.
\end{itemize}

\medskip\hrule\medskip

\section{Captain \textnormal{(Mercenary Companies)}}

The tactician and the muster: your archetype is your stock. [PROPOSAL C4: Troop rules in header.]

\subsection{The Creche Captain}
\textit{Cloned brothers --- expensive, capable, and people, whatever the Rolls decide.} Playstyle: elite small squad; troops that flank, hold, and grieve.
\begin{itemize}
    \item \textbf{Muster: Brothers} --- 1 --- L1 --- Troop ×2 --- Two Marchers: HP 3, Strike 3, Defend 1, Move 1 each. They use real names.
    \item \textbf{By the Numbers} --- 1 --- L2 --- Two of your troops acting together this surge: their attacks gain +1 die and share hits against one target.
    \item \textbf{The Line Holds} --- 1 --- L3 --- Stance --- Your troops within 3 tiles of you gain +1 Defend and cannot be routed or frightened.
    \item \textbf{Ninth of His Line} --- 2 --- L6 --- One-Shot --- One troop acts as a full character for the scene (your class dice, its body) --- and afterward, it asks you something.
\end{itemize}

\subsection{The Rite-Captain}
\textit{The leased dead --- tireless, fearless, contract-bound, and buried properly at term.} Playstyle: attrition; the wall that does not eat, sleep, or exceed the order.
\begin{itemize}
    \item \textbf{Muster: the Leased} --- 1 --- L1 --- Troop ×2 --- Two leased dead: HP 4, Strike 2, Defend 0, Move 1; immune to fear, poison, fatigue; \textit{never} exceed orders.
    \item \textbf{Stand and Take It} --- 1 --- L2 --- Stance --- Your dead within 3 tiles gain +1 Armor and hold ground against any push.
    \item \textbf{The Terms Are Clear} --- 1 --- L3 --- Reaction --- When a troop of yours would be turned, commanded, or corrupted by an outside power: it isn't. The lease is explicit.
    \item \textbf{Wages to the Heirs} --- 2 --- L6 --- One-Shot --- A destroyed troop rises once more until end of scene at full HP; at scene's end you personally close its lease, and the table sees you do it.
\end{itemize}

\subsection{The Verdant Captain}
\textit{The grown line --- patience, thorns, and fortifications that heal.} Playstyle: area denial; the map slowly becomes yours.
\begin{itemize}
    \item \textbf{Muster: Thorn-Soldiers} --- 1 --- L1 --- Troop ×2 --- Two grown: HP 3, Strike 2, Defend 1, Move 1; enemies that Strike them take 1 hit.
    \item \textbf{Root} --- 1 --- L2 --- A grown troop roots: Move 0, +2 Armor, +1 Strike, and its tile plus adjacent become difficult ground for enemies.
    \item \textbf{The Hedge Rises} --- 1 --- L4 --- Grow a 3-tile thorn wall (Defend 1, HP 4 per tile); it regrows 2 HP per round.
    \item \textbf{Third Line's Answer} --- 2 --- L7 --- One-Shot --- All your grown troops root as a free action and fight at +1 die for the scene; when you sound the recall, note carefully whether they come.
\end{itemize}

\subsection{The Chimeric Captain}
\textit{The beast-lines --- bred for the work, directed by the hand.} Playstyle: mobility and scouting; the pack sees first and arrives first.
\begin{itemize}
    \item \textbf{Muster: the Pack} --- 1 --- L1 --- Troop ×2 --- Two beasts: HP 2, Strike 2, Defend 2, Move 3; they flank by instinct (+1 die when paired on a target).
    \item \textbf{Loose the Wing} --- 1 --- L2 --- A flyer scouts: reveal everything within 10 tiles of a point you choose; the GM answers as the bird saw it.
    \item \textbf{Bring Them to Ground} --- 1 --- L3 --- A pack troop's hit also knocks the target down or holds it in place (your choice).
    \item \textbf{Bred for One Night} --- 2 --- L6 --- One-Shot --- Deploy the kennel's masterpiece for one scene: HP 6, Strike 4, Defend 2, Move 3, and a rider chosen at deployment (silent / climbing / carrying).
\end{itemize}

\subsection{The Free Captain}
\textit{Mortal colors --- a handful of volunteers, and the trust nothing made can buy.} Playstyle: leadership; no troops --- the party is the muster.
\begin{itemize}
    \item \textbf{You Heard the Captain} --- 1 --- L1 --- An ally acts immediately at --1 Momentum cost (minimum 0 for a Move).
    \item \textbf{Answer for Yours} --- 1 --- L2 --- Reaction --- When an ally within 5 tiles is hit, reduce the hits by 1; you flinch instead.
    \item \textbf{The Ledger} --- 1 --- L3 --- Once per session per settlement: name a contact from your years on the roads (a factor, a Warden, an old sergeant); the GM makes them real and reachable.
    \item \textbf{Hold Fast, Damn You} --- 2 --- L6 --- One-Shot --- Your side's Momentum this surge becomes 6 (do not roll). Afterward, everyone is spent: --1 die on everything for one round.
\end{itemize}

\medskip\hrule\medskip

\section{Bounty Hunter \textnormal{(Registries)}}

The Read, certified: one wake, followed to its end.

\subsection{The Wake-Reader}
\textit{Depth as doctrine --- the Hunter who knows the quarry better than the quarry does.} Playstyle: single-target supremacy; information becomes dice.
\begin{itemize}
    \item \textbf{Take the Name} --- 1 --- L1 --- Stance --- Designate your quarry (true name or established alias): +1 die on all rolls to find, follow, perceive, or fight it.
    \item \textbf{The Itch} --- 1 --- L1 --- Learn your quarry's bearing and rough distance; through walls, weather, and lies.
    \item \textbf{Dreamt It First} --- 1 --- L3 --- Reaction --- Your quarry acts against you or an ally: you anticipated it --- the affected creature gains +2 Defend or +2 dice on the save.
    \item \textbf{The Closing} --- 2 --- L6 --- One-Shot --- Against a quarry you've held for a full session: this scene, you act on its surges too (1 free action when its side activates). It dreams you back. Neither of you enjoys it.
\end{itemize}

\subsection{The Taker}
\textit{Terms: to answer, to hold --- the endgame specialist who delivers them breathing.} Playstyle: ambush and capture; the fight is over before it's a fight.
\begin{itemize}
    \item \textbf{From Concealment} --- 1 --- L1 --- Attack an enemy unaware of you: +2 dice and it can't apply Defend.
    \item \textbf{Shackle-Work} --- 1 --- L2 --- Against a downed, grappled, or surprised target: bind it (escape requires 2 successes vs your Knack).
    \item \textbf{Alive Means Alive} --- 1 --- L2 --- Stance --- Your attacks can always choose nonlethal, at no penalty, and you know exactly how close to the line you are.
    \item \textbf{Keeping the Appointment} --- 2 --- L6 --- One-Shot --- Confront your quarry alone-ish and speak: it must parley one full round --- and you gain, from its face, one true answer about who's behind the writ, if anyone is.
\end{itemize}

\subsection{The Veiled}
\textit{The counter-hunter --- warded, masked, and rude to things that swim the dreaming.} Playstyle: anti-psychic and defense; hunts what hunts by familiarity.
\begin{itemize}
    \item \textbf{Thin Your Name} --- 1 --- L1 --- Stance --- Supernatural attempts to locate, read, or target you specifically take --2 dice; mundane ones --1.
    \item \textbf{Feel the Read} --- 1 --- L2 --- You know when anything is tracking you or an adjacent ally by unnatural means, and its rough direction.
    \item \textbf{Turn the Thread} --- 1 --- L4 --- Reaction --- Something reads you: follow it back --- learn its nature and bearing, and it knows it was seen.
    \item \textbf{The Mask's Old Purpose} --- 2 --- L6 --- One-Shot --- For a scene, you and adjacent allies are \textit{unfamiliar}: creatures and effects that hunt by identity, name, or memory cannot hold you as targets. The Nol were right all along.
\end{itemize}

\medskip\hrule\medskip

\section{Shapeshifter \textnormal{(Frontier Wardens)}}

Forms come from your hearth --- the class's level-budget of shapes [PROPOSAL C6]; these pools are the traditions' techniques laid over them.

\subsection{The Pelt-Binder}
\textit{Few shapes, deeply held --- the woman who is a wolf, and has been its colleague thirty years.} Playstyle: committed transformation; the bound shape at full power.
\begin{itemize}
    \item \textbf{Don the Pelt} --- 1 --- L1 --- Shift into a bound form as part of any Move; while worn, that form's rolls gain +1 die.
    \item \textbf{The Skin Holds} --- 1 --- L2 --- Reaction --- Damage that would force you out of form doesn't; reduce it by 2 besides.
    \item \textbf{Vouched For} --- 1 --- L3 --- In your bound shape, Wild creatures of its kind and their kin read you as kin: passage, parley, or one favor, GM adjudicated.
    \item \textbf{Old Colleagues} --- 2 --- L6 --- One-Shot --- Your bound form fights beside you instead: it steps out (use the form's full statline as a Troop for the scene) while you remain yourself. Do not discuss this with the Enclaves.
\end{itemize}

\subsection{The Hearth-Walker}
\textit{Many light guests --- the cat for the town, the hawk for the sky, the horse for the road.} Playstyle: versatility; the right shape at the right instant.
\begin{itemize}
    \item \textbf{Quick Change} --- 1 --- L1 --- Shift as a free part of any action, once per surge.
    \item \textbf{Borrowed Instinct} --- 1 --- L2 --- In any form: use one of its senses or movement modes for a round while otherwise yourself.
    \item \textbf{The Household} --- 1 --- L3 --- Stance --- Each round, your first shift refunds its Momentum; your forms' passive traits linger 1 round after leaving them.
    \item \textbf{Many Things at Once} --- 2 --- L6 --- One-Shot --- For one round, layer two light forms: both statlines' best values, both traits, and the unmistakable look of something that has been many things and is not entirely sure what it currently is.
\end{itemize}

\subsection{The Face-Taker}
\textit{The town is also a landscape --- humanoid shapes, under the code.} Playstyle: infiltration and social control; the network's inside work.
\begin{itemize}
    \item \textbf{A Studied Face} --- 1 --- L1 --- Wear a face you've observed for a scene: Deceive to maintain gains +2 dice; the code requires you log it.
    \item \textbf{Crowd-Shape} --- 1 --- L2 --- In any settlement: become forgettably local --- pursuit and Perceive against you take --2 dice while you do nothing violent.
    \item \textbf{Wear the Office} --- 1 --- L4 --- Your face comes with its standing: doors, deference, and answers owed to whoever you appear to be, until tested against someone who truly knows them.
    \item \textbf{The Code's One Exception} --- 2 --- L7 --- One-Shot --- Wear the face of someone present and willing: for a scene, you are them to every sense and system --- writs, Reads, and recognition included. The network is told afterward. Always.
\end{itemize}

\medskip\hrule\medskip

\section{Table Notes}

\textbf{Archetype is silhouette.} A pool should be readable across a dark room: the Line-Holder plants, the Open Channel glows, the Rite-Captain's wall does not breathe. When building the full pools out from these samples, protect each archetype's verbs --- versatility belongs to the Book-Bound and the Hearth-Walker \textit{because} the Ward-Wright and the Pelt-Binder don't have it.

\textbf{The pools carry the lore.} Every archetype above is an institution's fingerprint, and several cards deliberately point at live setting questions --- the frame that chooses, the thorn line and its recall, the brother who asks, the dead man's lease closed on-screen. GMs should let those cards be doors: a Verdant Captain playing \textit{Third Line's Answer} has invited the Waking Made to their table, and the game intends exactly that.

\textbf{Sample, not census.} Three to four cards sketch each pool's grammar; full pools want 15--20 cards across levels 1--9, with duplicates-permitted marked per card. The costs, numbers, and level placements here are tuning stakes in the ground [PROPOSAL C7] --- the shapes are the design; the integers are negotiable.
